% RESUMO - PT
\begin{resumo}
  \vspace{-15pt}
  
  Este trabalho aborda a detecção e classificação de arritmias cardíacas em sinais de eletrocardiograma (ECG), utilizando técnicas de inteligência 
  artificial para apoiar o diagnóstico clínico e o monitoramento automatizado. O objetivo principal é desenvolver e comparar modelos de redes 
  neurais convolucionais unidimensionais (CNN1D) e bidimensionais (CNN2D) aplicados à classificação de batimentos cardíacos na base de dados 
  MIT-BIH Arrhythmia Database, considerando cinco classes: batimentos normais, três tipos de arritmias e uma classe de batimentos não classificados. 
  O processamento dos sinais envolveu etapas de pré-processamento, segmentação dos batimentos e organização dos conjuntos de treino e teste, 
  seguido do projeto e ajuste das arquiteturas convolucionais. O desempenho dos modelos foi avaliado por meio das métricas acurácia, sensibilidade, 
  especificidade, precisão e F1-score. A CNN1D obteve média geral de acurácia de 0.9744, sensibilidade de 0.9360, especificidade de 0.9840, 
  precisão de 0.9366 e F1-score de 0.9361, enquanto a CNN2D alcançou acurácia média de 0.9736, sensibilidade de 0.9340, especificidade de 0.9835, 
  precisão de 0.9365 e F1-score de 0.9341. Esses resultados indicam uma leve superioridade do modelo unidimensional, demonstrando o potencial das 
  arquiteturas propostas para aplicação em sistemas de engenharia voltados ao monitoramento cardíaco e ao suporte à decisão em ambiente clínico.

  Palavras-chave: \palavraschaveemlinha
\end{resumo}


% RESUMO - EN
\begin{resumo}[Abstract]
  \vspace{-15pt}
  
  \begin{otherlanguage*}{english}
  This work addresses the detection and classification of cardiac arrhythmias in electrocardiogram (ECG) signals, using artificial intelligence 
  techniques to support clinical diagnosis and automated monitoring. The main objective is to develop and compare one-dimensional (CNN1D) and 
  two-dimensional (CNN2D) convolutional neural network models applied to heartbeat classification on the MIT-BIH Arrhythmia Database, considering 
  five classes: normal beats, three types of arrhythmias, and a class of unclassified beats. Signal processing involved stages of preprocessing, 
  heartbeat segmentation, and organization of training and testing sets, followed by the design and tuning of the convolutional architectures. 
  The models' performance was evaluated using accuracy, sensitivity, specificity, precision, and F1-score metrics. The CNN1D obtained an overall 
  average accuracy of 0.9744, sensitivity of 0.9360, specificity of 0.9840, precision of 0.9366, and F1-score of 0.9361, while the CNN2D achieved 
  an average accuracy of 0.9736, sensitivity of 0.9340, specificity of 0.9835, precision of 0.9365, and F1-score of 0.9341. These results indicate 
  a slight superiority of the one-dimensional model, demonstrating the potential of the proposed architectures for application in engineering systems 
  aimed at cardiac monitoring and decision support in a clinical environment.
  
  Keywords: \inlinekeywords
\end{otherlanguage*}
\end{resumo}